\begin{proposition}\hypertarget{tag-01Q5}{}\label{fam:src:A.4.3}\label{fam:st:A.4.3}
\StatementRecord{fam:src:A.4.3}{proved}{01Q5}{\SeriesRef{fam:src:A.4.2}, \SeriesRef{fam:src:A.1.1}, \SeriesRef{fam:src:A.2.1}, \SeriesRef{fam:src:A.3.1}; \cite[Proposition 4.7]{ViaclovskyLog2026}; \cite[Lemmas 6.6--6.9 and Proposition 6.15]{Yang2026}}

The non-isotriviality argument of Viaclovsky, \cite[Proposition 4.7]{ViaclovskyLog2026}, follows for the present family from \ref{fam:src:A.4.2} and \ref{fam:src:A.1.1}. The cited countability argument proves local non-isotriviality and positive Kuranishi dimension; the computation in \ref{fam:src:A.4.2} additionally proves nonzero derivative at every admissible parameter. The arithmetic method of Yang, \cite[Lemmas 6.6--6.9 and Proposition 6.15]{Yang2026}, specializes to the same centralizer, finite-norm and collar-gluing operations after replacing his displayed matrices by $A_1,A_2$; those replacements are evaluated in \ref{fam:src:A.2.1} and \ref{fam:src:A.3.1}.
\end{proposition}


\begin{proof}
\begin{enumerate}[label=\arabic*.,leftmargin=*]
\item On a fixed ordinary fibre, the three last periods remain fixed and span $\mathbb C^2$, whereas the first changes by $\delta e$. A marked complex-linear identification fixing those three periods must be the identity, so the marked period map is nonconstant. The set of integral markings of one torus is countable. A holomorphic nonconstant map from a disc has uncountable image; hence the unmarked fibre family is locally non-isotrivial. \ref{fam:src:A.1.1} transfers this to total spaces. This is exactly the implication used in the cited Proposition 4.7, whose statement and proof we have read. A nonconstant classifying map gives positive dimension of its local target. It does not by itself determine the derivative at the base point, as the explicit holomorphic map $z\mapsto z^2$ has nonconstant image and derivative zero at zero. \ref{fam:src:A.4.2} computes the derivative here.

\item Yang's cited lemmas use an integral commuting lattice map, an affine fibre lift, normalization of the root base change, and a norm-lattice defect. We have exhibited the respective maps and norm equations for the present matrices. His Proposition 6.15 obtains $6/r$ from the first-column scaling and its order-six norm image. Here the first-column scaling is one and the two computed norm images impose parity, yielding two. The common proof operations transfer through these explicitly stated matrix substitutions; no conclusion about one family is inferred merely from the label of the other. His nonzero absolute Kodaira--Spencer conclusion is obtained in the present setting by the explicit computation and intrinsic-pencil transfer above.
\end{enumerate}
\end{proof}